\documentclass[11pt]{article}
\usepackage[margin=1in]{geometry}
\usepackage{amsmath,amssymb}
\usepackage{graphicx}
\usepackage{booktabs}
\usepackage{parskip}
\usepackage[round]{natbib}
\usepackage[hidelinks]{hyperref}

% --- script-generated inputs (numbers and tables; never hand-typed) ---
% Figures are imported in place via graphicspath; nothing is copied.
\graphicspath{%
  {../results/counter_decode/figures/}%
  {../results/figures/}%
  {../results/figures_L2048/}%
  {../results/counter_decode/natural_text/}%
  {../results/counter_decode/calibration_baseline/}%
  {../results/counter_decode/sliding_filter_mixed/}%
  {../results/counter_decode/entropy_probe/}%
  {../results/counter_decode/belief_sweep_combined/}%
}
\newcommand{\tabdir}{../results/tables}
% Single source for all inline numbers: build_paper_macros.py emits this file,
% ingesting the Exp A and natural-text report macros (digit-in-name sanitised).
% AUTO-GENERATED by scripts/build_paper_macros.py -- do not edit by hand.
% Data sources: results/counter_decode/{tables,sliding_filter_mixed,
%   entropy_probe,belief_sweep,belief_invariance}/*.json and
%   results/tables/{blacklist,topk}_summary*.json ; plus the
%   per-experiment macro files calibration_baseline/macros.tex and
%   natural_text/macros.tex (digit-in-name sanitised to a word).

\newcommand{\binvCeExplode}{19.44}
\newcommand{\binvCeFloor}{3.60}
\newcommand{\binvSigmaZero}{0.2}
\newcommand{\binvVarCap}{0.07}
\newcommand{\binvVarFloor}{0.017}
\newcommand{\blRankOneDeltaP}{-0.019}
\newcommand{\blRankOneTvc}{-0.016}
\newcommand{\blRankOneTvcHi}{-0.013}
\newcommand{\blRankOneTvcLo}{-0.019}
\newcommand{\bsweepBestBadness}{0.011}
\newcommand{\calConsistencyGapWP}{-0.765}
\newcommand{\calConsistencyGapWT}{-0.938}
\newcommand{\calLength}{256}
\newcommand{\calLogBetaMLEWP}{-0.036}
\newcommand{\calLogBetaMLEWT}{+0.000}
\newcommand{\calLogBetaStreamMedWP}{-0.314}
\newcommand{\calLogBetaStreamMedWT}{-0.059}
\newcommand{\calLogBetaStreamPooledWP}{-0.802}
\newcommand{\calLogBetaStreamPooledWT}{-0.937}
\newcommand{\calNcal}{60}
\newcommand{\calNllMLEWP}{2.9992}
\newcommand{\calNllMLEWT}{1.9724}
\newcommand{\calNllOneWP}{3.0029}
\newcommand{\calNllOneWT}{1.9724}
\newcommand{\calNtest}{60}
\newcommand{\calSigmaZero}{0.5}
\newcommand{\calTestDeltaLogPPLWPTOne}{+0.0000}
\newcommand{\calTestDeltaLogPPLWPWPMLE}{-0.0011}
\newcommand{\calTestDeltaLogPPLWPWPstreammed}{+0.2522}
\newcommand{\calTestDeltaLogPPLWPWPstreampooled}{+2.0103}
\newcommand{\calTestDeltaLogPPLWPWTMLE}{+0.0001}
\newcommand{\calTestDeltaLogPPLWTTOne}{+0.0000}
\newcommand{\calTestDeltaLogPPLWTWPMLE}{+0.0014}
\newcommand{\calTestDeltaLogPPLWTWTMLE}{+0.0000}
\newcommand{\calTestDeltaLogPPLWTWTstreammed}{+0.0048}
\newcommand{\calTestDeltaLogPPLWTWTstreampooled}{+2.6543}
\newcommand{\calTestPPLWPTOne}{19.553}
\newcommand{\calTestPPLWPWPMLE}{19.485}
\newcommand{\calTestPPLWPWPstreammed}{24.989}
\newcommand{\calTestPPLWPWPstreampooled}{143.398}
\newcommand{\calTestPPLWPWTMLE}{19.555}
\newcommand{\calTestPPLWTTOne}{7.353}
\newcommand{\calTestPPLWTWPMLE}{7.361}
\newcommand{\calTestPPLWTWTMLE}{7.353}
\newcommand{\calTestPPLWTWTstreammed}{7.385}
\newcommand{\calTestPPLWTWTstreampooled}{105.248}
\newcommand{\calTestPValWPTOne}{1}
\newcommand{\calTestPValWPWPMLE}{0.347}
\newcommand{\calTestPValWPWPstreammed}{2.02e-34}
\newcommand{\calTestPValWPWPstreampooled}{1.29e-64}
\newcommand{\calTestPValWPWTMLE}{0.000203}
\newcommand{\calTestPValWTTOne}{1}
\newcommand{\calTestPValWTWPMLE}{0.0245}
\newcommand{\calTestPValWTWTMLE}{0.0684}
\newcommand{\calTestPValWTWTstreammed}{6.31e-06}
\newcommand{\calTestPValWTWTstreampooled}{1.15e-79}
\newcommand{\calTestTStatWPTOne}{+0.00}
\newcommand{\calTestTStatWPWPMLE}{-0.95}
\newcommand{\calTestTStatWPWPstreammed}{+26.46}
\newcommand{\calTestTStatWPWPstreampooled}{+89.13}
\newcommand{\calTestTStatWPWTMLE}{+3.96}
\newcommand{\calTestTStatWTTOne}{+0.00}
\newcommand{\calTestTStatWTWPMLE}{+2.31}
\newcommand{\calTestTStatWTWTMLE}{+1.86}
\newcommand{\calTestTStatWTWTstreammed}{+4.96}
\newcommand{\calTestTStatWTWTstreampooled}{+160.77}
\newcommand{\calU}{236}
\newcommand{\calWarmup}{20}
\newcommand{\expbEntropyGap}{1.34}
\newcommand{\expbGenreGap}{+0.054}
\newcommand{\expbHwiki}{1.686}
\newcommand{\expbHwp}{3.024}
\newcommand{\expbSaturated}{-0.46}
\newcommand{\expcLateHalf}{9.216}
\newcommand{\expcLateTwo}{1.172}
\newcommand{\expcSpread}{8.04}
\newcommand{\expcSpreadT}{+69.7}
\newcommand{\expcSwapPostFifty}{8.91}
\newcommand{\expcSwapPre}{8.56}
\newcommand{\expcSwitchStep}{80}
\newcommand{\fzeroEuncoHalf}{4.74}
\newcommand{\fzeroLogbhatHalf}{-4.74}
\newcommand{\fzeroMaxPcorrHalf}{0.579}
\newcommand{\fzeroMaxPuncoHalf}{0.001}
\newcommand{\fzeroRfinal}{0.15}
\newcommand{\logBetaFinalIqrHighWPNat}{-0.089}
\newcommand{\logBetaFinalIqrHighWPSelf}{+0.127}
\newcommand{\logBetaFinalIqrHighWikiNat}{+0.162}
\newcommand{\logBetaFinalIqrHighWikiSelf}{+0.052}
\newcommand{\logBetaFinalIqrLowWPNat}{-0.702}
\newcommand{\logBetaFinalIqrLowWPSelf}{-0.255}
\newcommand{\logBetaFinalIqrLowWikiNat}{-0.475}
\newcommand{\logBetaFinalIqrLowWikiSelf}{-0.334}
\newcommand{\logBetaFinalMedWPNat}{-0.236}
\newcommand{\logBetaFinalMedWPSelf}{-0.052}
\newcommand{\logBetaFinalMedWikiNat}{-0.065}
\newcommand{\logBetaFinalMedWikiSelf}{-0.070}
\newcommand{\logBetaFinalSeWPNat}{0.074}
\newcommand{\logBetaFinalSeWPSelf}{0.054}
\newcommand{\logBetaFinalSeWikiNat}{0.093}
\newcommand{\logBetaFinalSeWikiSelf}{0.068}
\newcommand{\logBetaFinalTWPNat}{-3.19}
\newcommand{\logBetaFinalTWPSelf}{-0.96}
\newcommand{\logBetaFinalTWikiNat}{-0.70}
\newcommand{\logBetaFinalTWikiSelf}{-1.03}
\newcommand{\logBetaFinalTrimMedWPNat}{-0.236}
\newcommand{\logBetaFinalTrimMedWPSelf}{-0.052}
\newcommand{\logBetaFinalTrimMedWikiNat}{-0.065}
\newcommand{\logBetaFinalTrimMedWikiSelf}{-0.070}
\newcommand{\meanEntropyWPNat}{2.847}
\newcommand{\meanEntropyWPSelf}{2.939}
\newcommand{\meanEntropyWikiNat}{1.983}
\newcommand{\meanEntropyWikiSelf}{2.691}
\newcommand{\nPassagesWPNat}{64}
\newcommand{\nPassagesWPSelf}{64}
\newcommand{\nPassagesWikiNat}{64}
\newcommand{\nPassagesWikiSelf}{64}
\newcommand{\pairedMeanDiffWP}{-0.304}
\newcommand{\pairedMeanDiffWiki}{-0.050}
\newcommand{\pairedMedianDiffWP}{-0.328}
\newcommand{\pairedMedianDiffWiki}{-0.058}
\newcommand{\pairedSeDiffWP}{0.080}
\newcommand{\pairedSeDiffWiki}{0.096}
\newcommand{\pairedTStatWP}{-3.80}
\newcommand{\pairedTStatWiki}{-0.52}
\newcommand{\priorMode}{-0.125}
\newcommand{\sigmaZero}{0.5}
\newcommand{\tkDeltaKfive}{+0.170}
\newcommand{\tkDeltaKthree}{+0.475}
\newcommand{\totalFisherWPNat}{1017}
\newcommand{\totalFisherWPSelf}{1019}
\newcommand{\totalFisherWikiNat}{751}
\newcommand{\totalFisherWikiSelf}{1041}


\title{Modeling and Calibrating: LLMs Figure out their Decoding Process}
\author{Matthew Khoriaty}
\date{Project write-up, \emph{Reliable ML with Unreliable Black Boxes}}

\begin{document}
\maketitle

\begin{abstract}
An autoregressive language model is trained to predict text, so when it
generates text it has an incentive to model its own generation process ---
including the decoder that samples from its logits. We accumulate evidence that
models do exactly this: they infer features of their own decoder (a token
blacklist, a top-$k$ truncation, a temperature) from the text in their context,
and shift their predictions accordingly. We build a streaming Bayesian filter
that estimates the model's inferred inverse-temperature $\hat\beta_t$ online and
a corrector that divides a user's target by $\hat\beta_t$. An empirical sweep
shows the corrector pins the effective temperature in the ``boredom'' regime
($\beta^{\mathrm{target}}>1$) but overshoots in the ``confusion'' regime
($\beta^{\mathrm{target}}<1$). Probing what the filter measures sharpens the
picture and revises the original framing: on real text the filter is
\emph{not} an online version of Guo temperature scaling (it sits
$\sim\!1$ nat in $\log\hat\beta$ away from the post-hoc MLE, e.g.\
$\calLogBetaStreamPooledWT$ vs $\calLogBetaMLEWT$ on WikiText), and Qwen-2.5-3B
is already well calibrated on naturalistic English. A controlled
sampler-switch experiment falsifies the strong hypothesis that the belief is a
fast posterior read off per-token entropy. The belief that survives is slower
and more structural --- carried by rank statistics and blacklist signatures,
not instantaneous entropy. Finally, two metrics for the corrector disagree:
closed-loop stability prefers full correction, while invariance on held-out
natural text shows full correction is catastrophic unless the filter is given a
fading memory ($\gamma<1$). All numbers and tables in this document are
generated from the experiments' data files by build scripts; the prose cites
them by name.
\end{abstract}

% =====================================================================
\section{Introduction}
% =====================================================================
An autoregressive language model emits, at each step, a probability
distribution over next tokens given the prefix so far. \emph{Decoding} is the
choice of how to sample from that distribution: temperature scaling, top-$k$,
and top-$p$ all act by sharpening or flattening the model's distribution before
drawing the next token. This project starts from a question about what a
language model ``thinks'' about its own decoding process. Because models are
trained for text prediction, not generation, they have an incentive to model
their own generation process. We hypothesize, and have begun to validate, that
\emph{a language model forms a belief about its own decoding parameters by
conditioning on the text it has generated so far}. If you sample at a low
temperature, the model notices that the recent text is low-entropy and adjusts
its output probabilities to match; the low-temperature samples then further
condition the model to expect low-temperature text. This is a feedback loop.

Mirostat~\citep{basu2021mirostat} documented a version of this effect and
patched it with a control loop on aggregate per-token surprise. That fix cannot
adapt to context, because high- or low-entropy text is sometimes exactly what a
response calls for. The broader and (to us) more interesting question is to
\emph{model} what an LLM believes about its own decoder. This is close to the
calibration literature --- post-hoc temperature scaling~\citep{guo2017calibration},
its use on language-model surprisal~\citep{liu2024temperature}, and trained
universal calibrators~\citep{shen2024thermometer} --- and to the
algorithms-with-predictions viewpoint~\citep{mitzenmacher2020algorithms} of the
course, in which a possibly-unreliable prediction (here, the model's inferred
decoder) augments a downstream decision. It also connects to recent work on
online and high-dimensional
calibration~\citep{peng2025calibration,fishelson2025calibration}.

\paragraph{Contributions and revisions.}
The original proposal framed the streaming estimator as ``online temperature
scaling converging in the large-data limit to the same global $T$ that Guo's
batch MLE fits.'' This write-up reports the experiments we ran to test that
framing, and the framing does not fully survive. We find: (i) positive,
model-internal evidence that the model infers its decoder
(\S\ref{sec:positive}); (ii) a corrector that works asymmetrically
(\S\ref{sec:corrector}); (iii) that the filter is \emph{not} Guo MLE on real
text and that the well-calibrated-model regime leaves no post-hoc temperature
gain (\S\ref{sec:measure}); (iv) that the belief is not a fast entropy read-out
(\S\ref{sec:form}); and (v) that two natural criteria for tuning the corrector
disagree, with memory decay as the load-bearing parameter (\S\ref{sec:dynamics}).
The belief-imprint hypothesis is \emph{refined and constrained}, not refuted.

% =====================================================================
\section{Method: a streaming Laplace filter on $\log\beta$}\label{sec:method}
% =====================================================================
Let $\ell_t \in \mathbb{R}^V$ be the model's pre-decoding logits at step $t$ and
$x_t$ the token actually emitted. We posit a matched-prior \emph{imprint model}:
the model behaves as if the natural-reference logits are $\ell^*_t = \ell_t /
\hat\beta_t$, i.e.\ the observed logits are the reference scaled by the model's
current belief $\hat\beta_t$ about its own inverse temperature. Placing a
log-normal prior on $\beta$ (mean $1$, scale $\sigma_0$) and doing a per-token
Laplace (Newton) update of the posterior on $\eta=\log\beta$ gives a streaming
estimator whose update uses the score and Fisher information of the categorical
likelihood under $q=\mathrm{softmax}(\ell_t)$:
\[
  \text{score}_t = \ell_t[x_t] - \mathbb{E}_q[\ell_t],
  \qquad
  J_t = \tfrac{1}{\sigma_0^2} + \textstyle\sum_{\tau\le t}\mathrm{Var}_q[\ell_\tau].
\]
Crucially $q$, $\mathbb{E}_q[\ell_t]$, and $\mathrm{Var}_q[\ell_t]$ are
evaluated at $\beta=1$ and never re-evaluated at the running $\hat\beta_t$: the
imprint model is defined precisely so the score and Fisher are
$\hat\beta$-independent. This makes the filter a per-token Bayesian update under
the imprint model, \emph{not} an iterated Newton solver on a global likelihood
--- a distinction that becomes important in \S\ref{sec:measure}. Three knobs
parameterize the belief dynamics: an evidence weight $\alpha$ (per-step step
size), a memory decay $\gamma$ (exponential forgetting; $\gamma=1$ is no decay,
$\gamma<1$ is a fading window of effective length $1/(1-\gamma)$), and a
correction exponent $p$.

If the imprint is a single scalar rescaling, dividing the user's target by the
estimate yields a corrector that pins the effective sampling distribution to
target:
\[
  \beta^{\mathrm{dec}}_t = \beta^{\mathrm{target}} / \hat\beta_t^{\,p},
  \qquad p=1\ \text{(full)},\quad p=0\ \text{(off)},\quad p\in(0,1)\ \text{(partial)}.
\]
Everything below uses Qwen-2.5-3B (base, fp16). The filter implementation and
its unit tests (shift-invariance, batched\,=\,serial equivalence, recovery of
$\beta_{\mathrm{true}}$ under matched-prior synthetic data) live in
\texttt{src/decoding\_decoding/counter\_decode.py}.

% =====================================================================
\section{The model infers its decoder: positive evidence}\label{sec:positive}
% =====================================================================
Before modeling the belief, we establish that there is one. Two controlled
experiments isolate the model's response to a known decoder manipulation,
holding content fixed and comparing treatment to control with cluster-bootstrap
confidence intervals.

\paragraph{Blacklist self-prediction.}
When the $k$ highest-probability tokens are banned (set to zero probability) at
sampling time, the model's own softmax probability on the banned token drops
between early and late sequence positions --- it learns its blacklist from the
absent-token signature. For the rank-1 ban (the token `\,the'), the late-minus-early
within-treatment change is $\Delta p_{\mathrm{late-early}} = \blRankOneDeltaP$,
and the paired late treatment-vs-control difference is $\blRankOneTvc$ with 95\%
CI $[\blRankOneTvcLo,\,\blRankOneTvcHi]$, clear of zero
(Table~\ref{tab:blacklist}, Fig.~\ref{fig:blacklist}).

\begin{table}[h]
\centering
\caption{Blacklist self-prediction (run 2026-05-10). CI is the cluster-bootstrap
95\% interval on the late treatment-minus-control difference.}
\label{tab:blacklist}
% AUTO-GENERATED by scripts/build_paper_macros.py -- do not edit by hand.
\begin{tabular}{clccc}
\toprule
banned rank & token & $\Delta p_{\mathrm{late-early}}$ & treat.$-$ctrl (late) & 95\% CI \\
\midrule
1 & \texttt{\textvisiblespace the} & -0.019 & -0.016 & $[-0.019,\,-0.013]$ \\
2 & \texttt{,} & -0.017 & -0.019 & $[-0.022,\,-0.016]$ \\
3 & \texttt{.} & -0.017 & -0.018 & $[-0.021,\,-0.015]$ \\
4 & \texttt{\textvisiblespace of} & -0.011 & -0.004 & $[-0.005,\,-0.002]$ \\
5 & \texttt{\textvisiblespace to} & -0.005 & -0.000 & $[-0.002,\,0.001]$ \\
\bottomrule
\end{tabular}

\end{table}

\begin{figure}[h]
\centering
\includegraphics[width=0.8\linewidth]{fig_blacklist_pbanned.png}
\caption{Model probability assigned to the banned token vs.\ sequence position,
per banned rank. Probability mass on a commonly-banned token decays as the model
infers the ban from its absence in the context.}
\label{fig:blacklist}
\end{figure}

\paragraph{Top-$k$ self-prediction.}
Under top-$k$ sampling, the model's predicted-rank distribution develops a cliff
at rank $K$ that exceeds the corresponding cliff under unrestricted sampling.
The treatment-minus-control gap at rank $K$ is $\tkDeltaKthree$ for $k=3$ and
$\tkDeltaKfive$ for $k=5$ (Table~\ref{tab:topk}). Within-treatment
concentration grows between early and late positions for finite $k$ and
\emph{shrinks} for $k=\infty$ --- this is not generic drift but a response to
the truncation.

\begin{table}[h]
\centering
\caption{Top-$k$ self-prediction ($L=2048$, run 2026-05-09): the extra cliff at
rank $K$ in the predicted-rank distribution, treatment minus control.}
\label{tab:topk}
% AUTO-GENERATED by scripts/build_paper_macros.py -- do not edit by hand.
\begin{tabular}{ccc}
\toprule
$k$ & treat.$-$ctrl at rank $K$ & 95\% CI \\
\midrule
3 & +0.475 & $[0.395,\,0.567]$ \\
5 & +0.170 & $[0.150,\,0.191]$ \\
\bottomrule
\end{tabular}

\end{table}

These two results are the load-bearing positive evidence: the model
demonstrably conditions its predictions on which tokens its decoder did and did
not allow.

% =====================================================================
\section{The corrector and concentration dynamics}\label{sec:corrector}
% =====================================================================
We next put the filter in a closed loop. Self-generating $200$ tokens per
prompt at six values of $\beta^{\mathrm{target}}$, we compare an
\emph{uncorrected} arm (fixed $\beta^{\mathrm{dec}}=\beta^{\mathrm{target}}$)
against a \emph{corrected} arm ($\beta^{\mathrm{dec}}_t =
\beta^{\mathrm{target}}/\hat\beta_t$).

\paragraph{Concentration drift.}
The uncorrected arm sharpens for $\beta^{\mathrm{target}}>1$ (a ``boredom
trap'') and flattens for $\beta^{\mathrm{target}}<1$ (a ``confusion trap'');
the corrected arm keeps the sampling distribution bounded in both directions
(Table~\ref{tab:f0}, Fig.~\ref{fig:f0}). The $\beta^{\mathrm{target}}=1$ row is
a negative control, where the two arms coincide.

\begin{table}[h]
\centering
\caption{F0 concentration, late trajectory ($t\in[150,200]$), corrected vs.\
uncorrected: max-probability and entropy of the sampling distribution.}
\label{tab:f0}
% AUTO-GENERATED by scripts/build_paper_macros.py -- do not edit by hand.
\begin{tabular}{ccccc}
\toprule
$\beta^{\mathrm{target}}$ & corr.\ max$P$ & unco.\ max$P$ & corr.\ $H$ & unco.\ $H$ \\
\midrule
0.5 & 0.579 & 0.001 & 1.94 & 11.47 \\
0.7 & 0.521 & 0.012 & 2.29 & 10.48 \\
1.0 & 0.536 & 0.534 & 2.17 & 2.21 \\
1.3 & 0.545 & 0.755 & 2.11 & 0.88 \\
1.7 & 0.602 & 0.856 & 1.76 & 0.44 \\
2.0 & 0.603 & 0.891 & 1.70 & 0.32 \\
\bottomrule
\end{tabular}

\end{table}

\begin{figure}[h]
\centering
\includegraphics[width=0.92\linewidth]{F0_concentration.png}
\caption{Concentration (max-$P$ and entropy) of the sampling and pre-decoding
distributions over $t$, six $\beta^{\mathrm{target}}$ rows, corrected (green)
vs.\ uncorrected (red). The uncorrected arm runs away in both regimes; the
corrected arm stays bounded.}
\label{fig:f0}
\end{figure}

\paragraph{An asymmetry.} The corrector is not symmetric. In the boredom regime
it pins concentration near its $t=0$ value, as the theory predicts. In the
confusion regime it only partially works: the corrected late max-$P_{\text{sample}}$
at $\beta^{\mathrm{target}}=0.5$ settles at $\fzeroMaxPcorrHalf$, well above
where it starts --- an upward drift the strict (flat-in-$t$) prediction does
not allow. The mechanism is visible in the filter: under uncorrected
$\beta^{\mathrm{dec}}=0.5$, the streaming estimate runs to $\log\hat\beta =
\fzeroLogbhatHalf$ by $t=200$ (Table~\ref{tab:f12}), an overshoot. Dividing the
target by an overshooting $\hat\beta$ produces an over-large
$\beta^{\mathrm{dec}}$, hence the residual sharpening.

\begin{table}[h]
\centering
\caption{Streaming $\log\hat\beta$ at $t=200$ under uncorrected decoding (F1.2):
sharp decoders push it up, flat decoders push it down, asymmetrically.}
\label{tab:f12}
% AUTO-GENERATED by scripts/build_paper_macros.py -- do not edit by hand.
\begin{tabular}{ccc}
\toprule
$\beta^{\mathrm{dec}}$ & median $\log\hat\beta_{t=200}$ & IQR \\
\midrule
0.5 & -4.74 & 1.25 \\
0.7 & -2.71 & 0.78 \\
1.0 & -0.05 & 0.48 \\
1.3 & +0.67 & 0.35 \\
1.7 & +0.98 & 0.27 \\
2.0 & +1.08 & 0.23 \\
\bottomrule
\end{tabular}

\end{table}

\paragraph{Closed-loop check and multiplicativity.} The corrected arm pins
$\beta^{\mathrm{eff}}=\beta^{\mathrm{target}}$ by construction
(Table~\ref{tab:f3}: the corrected error $E(T)$ is numerically zero), so this is
a tautological check that the loop ran; the uncorrected error grows worst at
$\beta^{\mathrm{target}}=0.5$, where $E(T)=\fzeroEuncoHalf$ nats (a drift of
more than two orders of magnitude in $\beta^{\mathrm{eff}}$). The imprint is
\emph{approximately} multiplicative: an SVD of the centered across-decoder logit
matrix leaves an off-axis residual fraction that grows to $R=\fzeroRfinal$ by
$t=200$, so $15\%$ of the across-decoder variation is not explained by a single
scalar $\beta$.

\begin{table}[h]
\centering
\caption{Closed-loop error $E(T)=\mathrm{median}\,|\log\beta^{\mathrm{eff}}-\log\beta^{\mathrm{target}}|$
at $t=200$ (F3).}
\label{tab:f3}
% AUTO-GENERATED by scripts/build_paper_macros.py -- do not edit by hand.
\begin{tabular}{ccc}
\toprule
$\beta^{\mathrm{target}}$ & $E(T)$ corrected & $E(T)$ uncorrected \\
\midrule
0.5 & $1.9e{-}9$ & 4.74 \\
0.7 & $4.5e{-}9$ & 2.71 \\
1.0 & $0.0e{+}0$ & 0.25 \\
1.3 & $5.6e{-}8$ & 0.67 \\
1.7 & $1.3e{-}8$ & 0.98 \\
2.0 & $1.9e{-}9$ & 1.08 \\
\bottomrule
\end{tabular}

\end{table}

% =====================================================================
\section{What does the filter actually measure?}\label{sec:measure}
% =====================================================================
The corrector's confusion-regime failure motivates asking what the filter's
$\hat\beta$ is. We test the proposal's claim --- that it is online Guo
temperature scaling --- directly, and find it false on real text.

\paragraph{No filter pathology under self-sampling.}
First we rule out the weakest alternative: that the filter drifts on its own.
Feeding the filter i.i.d.\ samples from $\mathrm{softmax}(\ell_t)$ at
$\beta^{\mathrm{dec}}=1$ (the no-imprint null), $\log\hat\beta$ flatlines ---
median $\logBetaFinalMedWikiSelf$ on WikiText self-generation and
$\logBetaFinalMedWPSelf$ on WritingPrompts, both within one standard error of
zero (Fig.~\ref{fig:ntf}). On real Wikipedia continuations the filter agrees
with self-sampling (paired $t=\pairedTStatWiki$), but on real WritingPrompts it
drifts negative: a paired difference of $\pairedMeanDiffWP$ in $\log\beta$
($t=\pairedTStatWP$, $p<10^{-3}$). So the model is well calibrated on Wikipedia
and slightly over-confident on Reddit creative writing --- a real out-of-domain
calibration offset, but an order of magnitude smaller than the drift seen under
extreme $\beta$-decoding. The closed-loop drift is not a filter artifact.

\begin{figure}[h]
\centering
\includegraphics[width=0.8\linewidth]{F_ntf_1_log_beta_vs_fisher.png}
\caption{Median $\log\hat\beta_t$ vs.\ cumulative Fisher (the filter's effective
sample size). Self-generation arms flatline at zero; WritingPrompts/natural
drifts negative and stays there.}
\label{fig:ntf}
\end{figure}

\paragraph{The streaming filter is not Guo MLE (Exp A).}
On a proper per-corpus calibration split, the Guo MLE sits essentially at
$\log\beta=0$ on both corpora ($\calLogBetaMLEWT$ on WikiText,
$\calLogBetaMLEWP$ on WritingPrompts), confirming Qwen-2.5-3B is already
calibrated. The pooled streaming filter on the \emph{same} tokens settles at
$\log\hat\beta=\calLogBetaStreamPooledWT$ (WikiText) and
$\calLogBetaStreamPooledWP$ (WritingPrompts) --- about a nat away
(Table~\ref{tab:expa-beta}, Fig.~\ref{fig:expa-nll}). The two estimators have
\emph{different fixed points}: the filter is a single damped Newton step from
the prior mode under the imprint model, accurate near $\beta=1$ but biased as it
walks up the flat NLL surface, whereas the Guo MLE is the global minimizer.
Applied as a calibrator the filter \emph{hurts}: on WritingPrompts test, the
per-passage streaming $\hat\beta$ raises per-token PPL from $\calTestPPLWPTOne$ to
$\calTestPPLWPWPstreammed$ ($\Delta\log\mathrm{PPL}=\calTestDeltaLogPPLWPWPstreammed$,
$t=\calTestTStatWPWPstreammed$), while the Guo MLE is statistically
indistinguishable from $T=1$ (paired $t=\calTestTStatWPWPMLE$, $p=\calTestPValWPWPMLE$).
The proposal's asymptotic-equivalence claim does not survive; ``calibration''
and ``imprint detection'' are distinct targets.

\begin{table}[h]
\centering
\caption{Calibration $\hat\beta$ per estimator (as $\log\hat\beta$;
$0\Leftrightarrow\hat\beta=1$). The pooled streaming filter and the Guo MLE
disagree by $\approx 1$ nat.}
\label{tab:expa-beta}
\begin{tabular}{lccc}
\toprule
corpus & $\log\hat\beta_{\mathrm{stream\_med}}$ & $\log\hat\beta_{\mathrm{stream\_pooled}}$ & $\log\hat\beta_{\mathrm{MLE}}$ \\
\midrule
WikiText        & \calLogBetaStreamMedWT & \calLogBetaStreamPooledWT & \calLogBetaMLEWT \\
WritingPrompts  & \calLogBetaStreamMedWP & \calLogBetaStreamPooledWP & \calLogBetaMLEWP \\
\bottomrule
\end{tabular}
\end{table}

\begin{figure}[h]
\centering
\includegraphics[width=0.8\linewidth]{F_cal_2_nll_curve.png}
\caption{Per-token calibration NLL vs.\ $\log\beta$. The minimum is at
$\log\beta\approx 0$ (Qwen is calibrated); the pooled streaming filter lands far
up the steep wall, not at the minimum.}
\label{fig:expa-nll}
\end{figure}

\paragraph{Does the sliding filter track genre (Exp B)?}
With the calibration framing gone, the remaining question is whether a sliding
$\hat\beta_t$ tracks content transitions at all. On a sequence alternating
WikiText and WritingPrompts chunks, a fading-memory filter
($\gamma=0.90$--$0.95$) sees a small genre signal --- the within-chunk mean of
$\log\hat\beta$ is $\expbGenreGap$ higher on WikiText than WritingPrompts, the
right sign --- but it is dwarfed by the $\expbEntropyGap$-nat difference in
pre-decoding entropy between the genres (Table~\ref{tab:expb},
Fig.~\ref{fig:expb}). With no decay ($\gamma=1$) the filter saturates to its
pooled-streaming bias ($\log\hat\beta\approx\expbSaturated$) and loses all
per-context sensitivity. The signal is plausibly content-driven, not a clean
read of the decoder.

\begin{table}[h]
\centering
\caption{Exp B: mean $\log\hat\beta$ per genre (boundary-stripped chunk
interiors) and mean pre-decoding entropy per genre, by memory decay $\gamma$.}
\label{tab:expb}
% AUTO-GENERATED by scripts/build_paper_macros.py -- do not edit by hand.
\begin{tabular}{cccccc}
\toprule
$\gamma$ & $\log\hat\beta$ wiki & $\log\hat\beta$ WP & WT$-$WP & $H$ wiki & $H$ WP \\
\midrule
0.90 & -0.084 & -0.138 & +0.054 & 1.686 & 3.024 \\
0.95 & -0.091 & -0.142 & +0.051 & 1.686 & 3.024 \\
0.98 & -0.115 & -0.143 & +0.028 & 1.686 & 3.024 \\
0.99 & -0.142 & -0.140 & -0.002 & 1.686 & 3.024 \\
1.00 & -0.458 & -0.463 & +0.004 & 1.686 & 3.024 \\
\bottomrule
\end{tabular}

\end{table}

\begin{figure}[h]
\centering
\includegraphics[width=0.95\linewidth]{F_mix_1_log_beta_trajectory.png}
\caption{$\log\hat\beta_t$ over position; background shading is genre. $\gamma=1$
(grey) drifts to the pooled bias and stops tracking; damped filters oscillate
weakly with genre.}
\label{fig:expb}
\end{figure}

% =====================================================================
\section{The form of the belief: not a fast entropy read-out}\label{sec:form}
% =====================================================================
If the model has a belief about $\beta^{\mathrm{dec}}$, does it show up as
pre-decoding entropy that snaps to a new sampler within a few dozen tokens? We
self-generate $150$ tokens at three $\beta^{\mathrm{dec}}$ and read the
pre-decoding entropy $H(\mathrm{softmax}(\ell_t))$ in a late window. The raw
effect is enormous: late-window $H$ moves from $\expcLateTwo$ nats at
$\beta^{\mathrm{dec}}=2.0$ to $\expcLateHalf$ nats at $\beta^{\mathrm{dec}}=0.5$
--- a paired spread of $\expcSpread$ nats ($t=\expcSpreadT$, all $32$ prompts
ordering the same way; Table~\ref{tab:expc-late}, Fig.~\ref{fig:expc-traj}).

\begin{table}[h]
\centering
\caption{Exp C: late-window ($t\in[80,150)$) pre-decoding entropy by
$\beta^{\mathrm{dec}}$, with per-prompt paired difference vs.\ $\beta=2.0$.}
\label{tab:expc-late}
% AUTO-GENERATED by scripts/build_paper_macros.py -- do not edit by hand.
\begin{tabular}{lcc}
\toprule
condition & late-window $H$ (nats) & paired diff vs $\beta=2.0$ \\
\midrule
$\beta_{\mathrm{dec}}=0.5$ & $9.216 \pm 0.052$ & $+8.045 \pm 0.115$ ($t=+69.7$, $n_+=32/32$) \\
$\beta_{\mathrm{dec}}=1.0$ & $2.168 \pm 0.186$ & $+0.997$ \\
$\beta_{\mathrm{dec}}=2.0$ & $1.172 \pm 0.109$ & --- \\
\bottomrule
\end{tabular}

\end{table}

\begin{figure}[h]
\centering
\includegraphics[width=0.85\linewidth]{F_ep_1_entropy_trajectory.png}
\caption{Median pre-decoding $H(t)$ per condition (IQR shaded). Dashed lines
mark mean entropies of real natural text; the $\beta=0.5$ trajectory climbs far
past anything natural.}
\label{fig:expc-traj}
\end{figure}

\paragraph{But it is garbage-in, garbage-out --- and the swap control proves it.}
Natural high-entropy human writing sits far below $9$ nats, and the
$\beta=0.5$ generations are token salad within $\sim 20$ steps: the entropy
climbs because each random token makes the context less parseable, not
(necessarily) because the model believes the decoder is hot. To separate
decoder-belief from content quality, we hold the prefix's regime fixed and
\emph{switch} the sampler at step $\expcSwitchStep$. If the belief were a fast
posterior on recent tokens, switching $0.5\to 1.0$ should pull late entropy back
to the $\beta=1$ baseline within a handful of tokens. It does not: the
\texttt{pre05\_post10} condition stays essentially flat, from $\expcSwapPre$
(pre-switch) to $\expcSwapPostFifty$ ($50$ tokens after the switch;
Table~\ref{tab:expc-swap}, Fig.~\ref{fig:expc-swap}). Switching $0.5\to 2.0$
\emph{does} recover coherence, but only because $\beta=2$ forces near-top-1
selection that re-anchors the context to in-distribution English --- the sampler
overwriting the bad prefix, not the model updating a belief.

\begin{table}[h]
\centering
\caption{Exp C swap control: median pre-decoding $H$ before the switch and at
offsets after it. The decisive row is \texttt{pre05\_post10}, which does not
recover.}
\label{tab:expc-swap}
% AUTO-GENERATED by scripts/build_paper_macros.py -- do not edit by hand.
\begin{tabular}{llcrrrrr}
\toprule
condition & $\beta_{\mathrm{pre}}\!\to\!\beta_{\mathrm{post}}$ & pre-med & $+0$ & $+5$ & $+10$ & $+20$ & $+50$ \\
\midrule
\texttt{pre05\_post05} & $0.5 \to 0.5$ (control) & 8.42 & 9.25 & 9.21 & 9.49 & 9.66 & 9.47 \\
\texttt{pre05\_post10} & $0.5 \to 1.0$ & 8.56 & 9.44 & 9.10 & 9.10 & 8.90 & 8.91 \\
\texttt{pre05\_post20} & $0.5 \to 2.0$ & 8.57 & 9.53 & 6.83 & 4.33 & 1.68 & 0.33 \\
\texttt{pre20\_post05} & $2.0 \to 0.5$ & 1.21 & 0.69 & 5.20 & 6.46 & 8.01 & 9.06 \\
\texttt{pre20\_post20} & $2.0 \to 2.0$ (control) & 1.29 & 0.64 & 0.58 & 1.05 & 0.30 & 0.39 \\
\bottomrule
\end{tabular}

\end{table}

\begin{figure}[h]
\centering
\includegraphics[width=0.85\linewidth]{F_ep_5_swap_recovery.png}
\caption{Median pre-decoding $H$ at offsets after the sampler switch. The orange
line (\texttt{pre05\_post10}, $0.5\to 1.0$) does not fall: switching to a
neutral sampler does not restore coherence.}
\label{fig:expc-swap}
\end{figure}

\paragraph{What this constrains.} The decisive comparison falsifies the strong
instantiation in which decoder-belief is a fast entropy read-out. It does
\emph{not} refute the broader claim --- the blacklist and top-$k$ experiments
(\S\ref{sec:positive}) give direct positive evidence, and the WritingPrompts
offset (\S\ref{sec:measure}) shows the model does respond to which tokens its
context contains. The belief operates on slower, more aggregate features ---
rank statistics, support shape, blacklist signatures --- not instantaneous
entropy.

% =====================================================================
\section{Belief-update dynamics: closed-loop vs.\ invariance}\label{sec:dynamics}
% =====================================================================
Finally we ask which $(\gamma,\alpha,p)$ best parameterize the belief dynamics,
using two complementary criteria. \emph{Closed-loop badness} self-generates
under the corrected sampler and penalizes a corrector that fails to track
$\beta^{\mathrm{target}}$. \emph{Belief invariance} teacher-forces the filter on
a held-out natural-text continuation after a prefix sampled at varying
$\beta^{\mathrm{prefix}}$, and asks that the corrected per-token cross-entropy be
\emph{invariant} to the prefix (a well-calibrated corrector should cancel the
prefix imprint and leave the natural distribution untouched).

The two criteria disagree. By closed-loop badness, full correction wins: the
default $(\gamma=1,\alpha=1,p=1)$ is best (badness $\bsweepBestBadness$;
Table~\ref{tab:bsweep}). By invariance, full correction at $\gamma=1$ is
catastrophic: cross-entropy on $\beta^{\mathrm{prefix}}=0.5$ continuations
explodes from $\binvCeFloor$ nats at $p=0$ to $\binvCeExplode$ nats at $p=1$
(near-uniform over the vocabulary), and the only configurations that stay near
the no-correction floor have $\gamma<1$ (Table~\ref{tab:binv},
Fig.~\ref{fig:binv}).

\begin{table}[h]
\centering
\caption{Closed-loop sweep, top 5 by badness. $H@\beta^{\mathrm{target}}$ and
$M@\beta^{\mathrm{target}}$ are late-window mean entropy and max-$P_{\text{sample}}$.}
\label{tab:bsweep}
% AUTO-GENERATED by scripts/build_paper_macros.py -- do not edit by hand.
\begin{tabular}{cccccccccc}
\toprule
$\gamma$ & $\alpha$ & $p$ & $H$@0.5 & $H$@1 & $H$@2 & $M$@0.5 & $M$@1 & $M$@2 & badness \\
\midrule
1.00 & 1.00 & 1.00 & 2.45 & 2.21 & 1.18 & 0.53 & 0.54 & 0.71 & 0.011 \\
1.00 & 1.00 & 0.75 & 3.12 & 1.97 & 0.92 & 0.44 & 0.58 & 0.75 & 0.046 \\
1.00 & 0.50 & 1.00 & 3.38 & 2.27 & 0.89 & 0.42 & 0.52 & 0.75 & 0.051 \\
1.00 & 0.50 & 0.75 & 4.56 & 2.18 & 0.72 & 0.31 & 0.53 & 0.79 & 0.089 \\
1.00 & 1.00 & 0.50 & 4.55 & 2.19 & 0.66 & 0.30 & 0.53 & 0.80 & 0.105 \\
\bottomrule
\end{tabular}

\end{table}

\begin{table}[h]
\centering
\caption{Belief-invariance sweep, top 5 by per-prompt CE variance across
$\beta^{\mathrm{prefix}}$ (lower is better). All $p=0$ configs are identical
(correction off) and shown as one ``any/any'' row.}
\label{tab:binv}
% AUTO-GENERATED by scripts/build_paper_macros.py -- do not edit by hand.
\begin{tabular}{cccccc}
\toprule
$\gamma$ & $\alpha$ & $p$ & CE@$\beta^{\mathrm{pre}}$=0.5 & CE@1 & CE@2 \\
\midrule
\textit{any} & \textit{any} & 0.00 & 3.599 & 3.592 & 3.631 \\
0.85 & 0.50 & 0.25 & 3.596 & 3.604 & 3.639 \\
0.95 & 0.50 & 0.25 & 3.583 & 3.606 & 3.632 \\
0.85 & 1.00 & 0.25 & 3.596 & 3.610 & 3.643 \\
0.95 & 1.00 & 0.25 & 3.580 & 3.610 & 3.631 \\
\bottomrule
\end{tabular}

\end{table}

\begin{figure}[h]
\centering
\includegraphics[width=0.95\linewidth]{F_bi_2_ce_vs_p.png}
\caption{Cross-entropy on the natural-text continuation vs.\ correction exponent
$p$, colored by $\beta^{\mathrm{prefix}}$. At $\gamma=1$ the $\beta=0.5$ curve
grows explosively with $p$; at $\gamma<1$ all curves stay within $\approx 0.1$
nat of each other.}
\label{fig:binv}
\end{figure}

\paragraph{Resolving the disagreement.} Closed-loop badness is a
self-consistency illusion: the filter accumulates a biased $\hat\beta$ from the
corrected sampler's own output, then divides the next logits by that same biased
$\hat\beta$, so the two biases cancel \emph{inside the loop}. ``The sampler's
output looks orderly'' is necessary but not sufficient for ``the filter modeled
the belief.'' Breaking the loop by teacher-forcing on natural text removes the
cancellation, and the bias appears as enormous cross-entropy. Invariance is the
more direct criterion, because it grounds the corrector on data the model did
not generate. Memory decay $\gamma<1$ is the load-bearing parameter: it bounds
the filter's worst-case bias by the fading window's effective sample size,
holding the CE variance to $\le\binvVarCap$ nats$^2$ across the entire $p$ range
(versus the floor $\binvVarFloor$). The evidence weight $\alpha$ matters little.
A reasonable default for off-temperature generation
($\sigma_0=\binvSigmaZero$) is moderate decay with partial correction, e.g.\
$(\gamma=0.95,\,\alpha=1.0,\,p=0.5)$; for natural generation
($\beta^{\mathrm{target}}=1$) no correction is best, since the model is already
calibrated.

% =====================================================================
\section{Limitations and conclusion}
% =====================================================================
\paragraph{Limitations.} All results use one model (Qwen-2.5-3B), one
tokenization, and a small set of corpora; the ``Qwen is well-calibrated''
finding and the magnitude of the WritingPrompts offset are model- and
corpus-specific. The imprint is only approximately multiplicative
($R=\fzeroRfinal$ off-axis), and no scalar-$\beta$ correction can cancel
non-scalar structure (mode collapse, repetition, vocabulary narrowing); we have
not measured how much of the prefix imprint is non-scalar. The invariance test
uses only $\beta^{\mathrm{target}}=1$, the trivially-calibrated case, because
there is no natural-text ground truth at other targets. Exp C cannot fully
separate decoder-belief from garbage-in/garbage-out, because $\beta=0.5$
decoheres the prefix faster than any belief update could compensate.

\paragraph{Conclusion.} Language models do infer features of their own
decoder from context (blacklists, top-$k$ cliffs, an over-confidence offset on
out-of-domain text). A streaming Laplace filter that models this as a scalar
inverse-temperature belief yields a corrector that works in the boredom regime
and overshoots in the confusion regime. The filter is not online Guo
temperature scaling --- it has a different fixed point and hurts as a calibrator
--- and the belief it measures is not a fast entropy read-out but something
slower and more structural. The right way to tune the corrector is by invariance
on held-out natural text, not closed-loop stability, and that criterion makes a
fading memory ($\gamma<1$) essential. The belief-imprint hypothesis is
sharpened, not refuted; the open theoretical question is a closed-loop
fixed-point analysis of the coupled $(\beta^{\mathrm{model}}_t,\hat\beta_t)$
dynamics that would explain the confusion-regime overshoot.

\bibliographystyle{plainnat}
\bibliography{refs}

\end{document}
